\section{Introduction}
\label{sec:introduction}

Static analysis tools are routinely deployed in secure software engineering pipelines precisely because they are scalable and conservative: they sacrifice specificity for recall, flagging every potential vulnerability that cannot be ruled out by over-approximate reasoning. In industrial practice, this conservatism produces false-positive (FP) rates between 40\% and 90\%~\cite{johnson2013,christakis2016}, with manual triage requiring 10--20 minutes per alert~\cite{tencent2023}. At enterprise scale, the cumulative review burden is a material engineering cost that limits the practical impact of static analysis.

Recent work demonstrates that large language models (LLMs) can reduce this burden by conditioning on structured analyzer evidence---code slices, data-flow traces, and path constraints---to distinguish real vulnerabilities from analysis artifacts~\cite{llm4fpm2024,llm4pfa2024,buglens2024,adataint2023}. However, two deployment barriers persist. The closest prior system, IRIS~\cite{li2025iris}, uses LLM-assisted specification inference to validate CodeQL findings but requires generating per-vulnerability formal specifications and relies on frontier-model reasoning without principled abstention. First, state-of-the-art triage systems rely on computationally expensive frontier models; replacing them with fast, cost-effective ``lite'' models (e.g., Gemini Flash Lite, Claude Haiku, GPT nano) reveals a calibration collapse: our multi-model evaluation on 3,000+ Juliet alerts spanning 87 CWE classes finds zero-shot precision ceilings of 42--60\% and expected calibration errors (ECE) exceeding 0.35, confirming that raw model confidence does not reliably track correctness in the lite-model regime. Second, existing LLM triage systems force binary classification without formalizing \emph{abstention} as a first-class action with auditable guarantees---an omission that is unsafe when missed vulnerabilities carry asymmetric costs.

We address both barriers with \textbf{EVICT} (Evidence-Conditioned Investigation with Calibrated Triage). EVICT casts alert triage as a \emph{cost-sensitive selective prediction} problem~\cite{elyaniv2010,geifman2017}: the system learns both a predictor $h$ and a rejector $g$, where $g(x) = 0$ explicitly defers uncertain alerts to human review, and $g(x) = 1$ accepts only those for which calibrated evidence supports a reliable decision. This framing converts the accept/abstain threshold from an arbitrary hyperparameter into the solution of a finite-sample optimization problem, ties automation to formal risk control, and creates a principled trigger for targeted symbolic verification.

\paragraph{Contributions.}

\textbf{(1) Formalization.} We cast alert triage as a cost-sensitive selective prediction problem with asymmetric costs, derive a finite-sample excess-risk bound for threshold selection with label-specific convergence rates (Theorem~\ref{thm:risk}), and prove distribution-free conformal validity for the pre-escalation selective predictor (Theorem~\ref{thm:conformal}). A label-conditional extension with exact threshold optimization under asymmetric costs is established in Theorems~\ref{thm:label_conditional} and~\ref{thm:exact_opt}.

\textbf{(2) Evidence-conditioned reasoning with calibrated confidence.} EVICT extracts a structured EvidencePack (code slice, data-flow summary, path constraints) from SARIF-format analyzer output and prompts a lite LLM with a fixed four-step reasoning schema. Confidence is estimated via logit margins or self-consistency vote shares ($k = 5$), then calibrated using \textbf{Split-Conformal Prediction} (\textbf{Algorithm~\ref{alg:conformal}}). An alert is accepted as TP or FP only when the conformal prediction set is a singleton; otherwise the system abstains.

\textbf{(3) Conditional neuro-symbolic escalation.} When the calibrated stage abstains, EVICT invokes a symbolic escalator---\textbf{Java PathFinder (JPF)} for Java---on deferred alerts. Symbolic \texttt{UNKNOWN} results and timeouts (8.1\% empirical rate) are strictly treated as continued abstentions rather than forced classifications, preserving the safety semantics of the pre-escalation stage.

\paragraph{Empirical contributions.} We contribute three complementary empirical studies. A \emph{multi-model zero-shot baseline study} evaluates five lite LLMs across 87+ CWE classes, finding precision ceilings of 42--60\% and ECE $> 0.35$ throughout. A \emph{multi-model conformal calibration study} applies the full EVICT pipeline to six models spanning three training paradigms (reasoning, instruction-tuned, code-specialized) on CWE-89/78/190 with $k=5$ self-consistency sampling and 5-fold cross-validation. This study reveals a previously unreported finding: \textbf{the quality of the self-consistency confidence signal depends more on the model's training objective than on model size}. Reasoning models (R1-Distill-14B) produce diverse vote-share (36\% unanimous, ECE 0.335), while instruction-tuned and code-specialized models produce degenerate confidence (85--98\% unanimous, ECE 0.51--0.62), making conformal calibration effective only for the former. A \emph{real-world validation} on CWE-Bench-Java (549 CodeQL alerts, 45 projects) confirms the TP-bias pattern and reveals a code-specialization paradox: code models achieve the best FP rejection (TN=184) despite the worst confidence quality. We further project that frontier reasoning models (GPT-4.5, Claude Sonnet 4) will be the first model class where conformal calibration produces practically significant precision gains ($+3$--$10$pp).

Section~\ref{sec:background} surveys related work. Sections~\ref{sec:methodology} and~\ref{sec:theory} specify the method and its formal guarantees. Section~\ref{sec:preliminary} reports the two empirical studies. The remaining sections describe the planned full evaluation, discuss limitations, and address broader impact.
